\chapter{Praktische Realisierung}
\abschnittsautor{B. Beispiel}

Das Kernkapitel: Was habt ihr gebaut, welche Entscheidungen habt ihr dabei
getroffen und warum? Beschreibt Entscheidungen mit ihren Alternativen, nicht
nur das Endergebnis.

\section{Umsetzung Teilbereich A}

Kurze Codeausschnitte bindet ihr als Listing ein -- nicht als Screenshot:

\begin{lstlisting}[language=Python, caption={Beispiel für ein Quellcode-Listing}, label=code:beispiel]
def berechne_gesamtpreis(positionen):
    """Summiert die Positionen eines Warenkorbs."""
    return sum(p.menge * p.einzelpreis for p in positionen)
\end{lstlisting}

Auf Listings verweist ihr mit \verb|\ref|: siehe Quellcode~\ref{code:beispiel}.
Einzelne Bezeichner im Fließtext setzt ihr mit \verb|\texttt|, etwa
\texttt{berechne\_gesamtpreis}.

\subsection{Ein Detail der Umsetzung}

Beschreibt hier ein konkretes Problem und wie ihr es gelöst habt.

\section{Umsetzung Teilbereich B}
\abschnittsautor{V. Vorlage}

Jedes Teammitglied dokumentiert seinen eigenen Teil. Den Verfasser stellt ihr
direkt nach \verb|\section| mit \verb|\abschnittsautor| um -- das geht pro
Unterkapitel, nicht nur pro Kapitel. Teilen sich zwei Abschnitte eine Seite,
nennt die Kopfzeile beide.
